% Companion Version of the Monograph
% Contains Sections 1-3, 12, and Appendices (Notation, Worked Example)

\documentclass[11pt]{amsart}
\usepackage{amsmath,amssymb,amsfonts,amsthm}
\usepackage{mathtools}
\usepackage{mathrsfs} % For \mathscr command
\usepackage[T1]{fontenc}
\usepackage{lmodern}
\usepackage[final]{microtype} % Enhanced microtype settings
\usepackage{graphicx}
\usepackage{array} % Added for table column formatting
\usepackage{booktabs}
\usepackage{longtable} % For multi-page tables
\usepackage{enumitem}
\setlist{nosep} % Compact lists for amsart compatibility
\usepackage{tikz}
\usetikzlibrary{arrows.meta,calc,positioning,shapes}
\usepackage{tikz-cd}
\usepackage{hyperref}
\usepackage[capitalize]{cleveref}
\usepackage{framed} % For boxed theorems

% Help LaTeX avoid overfull boxes globally
% Settings optimized for mathematical text with long equations
\setlength{\emergencystretch}{3em}
\tolerance=9999
\hbadness=9999
\vbadness=9999
\hfuzz=120pt % Suppress overfull box warnings (typical for math papers)

% Print-friendly hyperlinks (black for print submission)
\hypersetup{
    colorlinks=true,
    linkcolor=black,
    citecolor=black,
    urlcolor=black
}

\theoremstyle{plain}
\newtheorem{theorem}{Theorem}[section]
\newtheorem*{theorem*}{Theorem}
\newtheorem{lemma}[theorem]{Lemma}
\newtheorem{proposition}[theorem]{Proposition}
\newtheorem{corollary}[theorem]{Corollary}
\newtheorem{claim}[theorem]{Claim}

\theoremstyle{definition}
\newtheorem{definition}[theorem]{Definition}
\newtheorem{assumption}[theorem]{Assumption}

\theoremstyle{remark}
\newtheorem{remark}[theorem]{Remark}
\newtheorem{example}[theorem]{Example}

\theoremstyle{plain}
\newtheorem{conjecture}[theorem]{Conjecture}
\newtheorem{openproblem}[theorem]{Open Problem}

\newtheorem{maintheorem}{Theorem}
\renewcommand{\themaintheorem}{\Alph{maintheorem}}

\numberwithin{equation}{section}

% MOTS: Marginally Outer Trapped Surface - a codimension-2 surface with vanishing
% outward null expansion theta^+ = H + tr_Sigma(k) = 0
\newcommand{\MOTS}{\text{MOTS}}
\renewcommand{\Cap}{\text{Cap}}
\newcommand{\tr}{\operatorname{tr}}
\newcommand{\Tr}{\operatorname{Tr}}
\newcommand{\Wkp}{W^{1,p}_{\text{loc}}}
\newcommand{\Hone}{H^1_{\text{loc}}}
\newcommand{\Eigen}{\lambda_1}
% Notation disambiguation for alpha:
% - \alphaH = H\"older exponent in C^{k,alpha} regularity (always in (0,1))
% - \alphaInd = indicial root at bubble tips (phi ~ r^alphaInd)
\newcommand{\alphaH}{\alpha_{\mathrm{H}}}
\newcommand{\alphaInd}{\alpha_{\mathrm{ind}}}
\newcommand{\Lap}{\Delta}
\newcommand{\Vol}{\operatorname{Vol}}
\newcommand{\Area}{\operatorname{Area}}
\newcommand{\ADM}{\mathrm{ADM}}
\newcommand{\R}{\mathbb{R}}
\newcommand{\geps}{g_{\epsilon}}
\newcommand{\hatgeps}{\hat{g}_{\epsilon}}
\newcommand{\Met}{\mathcal{M}}
\newcommand{\JOp}{\mathcal{J}}
\newcommand{\LOp}{\mathcal{L}}
\newcommand{\Jump}[1]{[\![ #1 ]\!]}
\newcommand{\Weight}[2]{W^{#1, p}_{#2}}
\newcommand{\Holder}[2]{C^{#1, \alphaH}_{#2}}
\newcommand{\Norm}[2]{\|#1\|_{#2}}
\newcommand{\EdgeSpace}[2]{H^{#1}_{0,#2}}
\newcommand{\Ind}{\mathrm{Ind}}
\newcommand{\ind}{\mathrm{ind}}
\newcommand{\coker}{\mathrm{coker}}
\newcommand{\dist}{\mathrm{dist}}
\newcommand{\Spec}{\mathrm{Spec}}
\newcommand{\Harm}{\mathcal{H}}
\newcommand{\Energy}{\mathcal{E}}
\newcommand{\bM}{\overline{M}}
\newcommand{\bg}{\overline{g}}
\newcommand{\tM}{\widetilde{M}}
\newcommand{\tg}{\widetilde{g}}
\newcommand{\Rg}{R_{\overline{g}}}
\newcommand{\Rtg}{R_{\widetilde{g}}}
\newcommand{\dV}{\,dV}
\newcommand{\dVol}{\,d\text{Vol}}
\newcommand{\dsigma}{\,d\sigma}
\newcommand{\Scal}{\mathrm{R}}
\newcommand{\Ric}{\mathrm{Ric}}
\newcommand{\Div}{\mathrm{div}}
\newcommand{\supp}{\mathrm{supp}}
\newcommand{\diam}{\mathrm{diam}}
\newcommand{\Crit}{\mathcal{C}}
\newcommand{\llbracket}{[\![}
\newcommand{\rrbracket}{]\!]}
\DeclareMathOperator*{\argmax}{argmax}
\DeclareMathOperator{\sgn}{sgn}
\newcommand{\divv}{\mathrm{div}}
\newcommand{\fint}{\mathchoice{\,-\!\!\!\!\!\int}{\,-\!\!\!\!\int}{\,-\!\!\!\int}{\,-\!\!\!\int}}

\title[Spacetime Penrose Inequality (Companion)]{The Spacetime Penrose Inequality: Conditional Results (Companion Version)}
\author{Da Xu}
\address{China Mobile Research Institute, Beijing, China}
\email{xudayj@chinamobile.com}
\date{December 2025}

\subjclass[2020]{Primary 83C57; Secondary 53C21, 35Q75, 35J60, 49Q22}
\keywords{Penrose inequality, marginally outer trapped surface, Jang equation, cosmic censorship, dominant energy condition}

\begin{document}

\begin{abstract}
This is a \textbf{short companion version} (approx.\ 30 pp) of the full monograph "The Spacetime Penrose Inequality: Conditional Results for Stable MOTS and General Trapped Surfaces". It contains the Introduction, Overview, and the Consolidated Master Proof (Section 12), omitting the heavy technical machinery of the intermediate sections. It is designed to provide a clear, audit-oriented roadmap of the conditional proof, highlighting the role of the favorable jump condition, the KKT-to-distributional-jump interface, and the reliance on Weak Cosmic Censorship or the "outermost maximizer" hypothesis for the general case.
\end{abstract}

\maketitle

\setcounter{tocdepth}{1}
\tableofcontents

\part{Introduction and Overview}
\input{sec_01_introduction}
\input{sec_02_the_penrose_conjecture}
\input{sec_03_overview}

\part{Consolidated Proof}
\input{sec_12_complete_rigorous_proof_consolidated_statement}

\part{Appendices}
\appendix
\section{Index of Notation}\label{sec:Notation}

To assist the reader, we summarize the principal symbols, spaces, and functionals used throughout the proof.

\begin{remark}[Notation Conventions]
\textbf{Jang metric:} We use $\bg$ (equivalently written as $\bar{g}$ or $\overline{g}$) for the Jang metric. The macro $\backslash$\texttt{bg} expands to $\overline{g}$.

\textbf{Weight parameters:} On cylindrical ends, $\beta$ denotes the exponential weight in Lockhart--McOwen spaces. On asymptotically flat ends, $\delta$ denotes the polynomial weight. The specific choice $\beta \in (-1, 0)$ for marginally stable MOTS avoids indicial roots.

\textbf{Eigenvalue indexing:} We use \textbf{1-indexing} for eigenvalues of the stability operator $L_\Sigma$. Thus $\lambda_1$ denotes the \textbf{principal (smallest)} eigenvalue, and a stable MOTS satisfies $\lambda_1 \ge 0$. A marginally stable MOTS has $\lambda_1 = 0$, while a strictly stable MOTS has $\lambda_1 > 0$.
\end{remark}

\begin{table}[ht]
\centering
\caption{Metrics, manifolds, and domains.}
\begin{tabular}{l l l}
\hline
\textbf{Symbol} & \textbf{Description} & \textbf{Regularity} \\
\hline
$(M, g, k)$ & Initial data set & Smooth ($C^\infty$) \\
$(\bM, \bg)$ & Jang manifold (graph of $f$) & Lipschitz ($C^{0,1}$) at $\Sigma$ \\
$(\tM, \tg)$ & Conformal deformation ($\tg = \phi^4 \bg$) & $C^0$ at tips $p_k$, Lipschitz at $\Sigma$ \\
$(\tM, \geps)$ & Smoothed manifold (Miao--Piubello) & Smooth ($C^\infty$) \\
$\Sigma$ & Outermost MOTS (horizon) & Smooth embedded surface \\
$\{p_k\}$ & Compactified Jang bubbles & Conical singularities \\
$\mathcal{E}_{cyl}$ & Cylindrical end over $\Sigma$ & $[0,\infty) \times \Sigma$ \\
$\mathcal{E}_{AF}$ & Asymptotically flat end & $M \setminus K$ for compact $K$ \\
\hline
\end{tabular}
\end{table}

\begin{table}[ht]
\centering
\caption{Function spaces and operators.}
\begin{tabular}{l l}
\hline
\textbf{Symbol} & \textbf{Description} \\
\hline
$W^{k,p}_{\delta,\beta}(\bM)$ & Weighted Sobolev space (Lockhart--McOwen) \\
$\Hone$ & $H^1_{\mathrm{loc}}$, locally $H^1$ functions \\
$\Wkp$ & $W^{1,p}_{\mathrm{loc}}$, locally $W^{1,p}$ functions \\
$L_\Sigma$ & Stability operator of MOTS: $L_\Sigma = -\Delta_\Sigma + \tfrac{1}{2}R_\Sigma - |A|^2 - \Ric(\nu,\nu)$ \\
$\Delta_p$ & $p$-Laplacian: $\Delta_p u = \Div(|\nabla u|^{p-2} \nabla u)$ \\
$\mathcal{L}_\phi$ & Lichnerowicz operator: $\mathcal{L}_\phi = -8\Delta + R$ \\
\hline
\end{tabular}
\end{table}

\begin{table}[ht]
\centering
\caption{Key functionals and scalar quantities.}
\begin{tabular}{l l}
\hline
\textbf{Symbol} & \textbf{Description} \\
\hline
$M_{\ADM}(g)$ & ADM mass of metric $g$ \\
$A(\Sigma)$ & Area of surface $\Sigma$ \\
$\mathcal{M}_p(t)$ & AMO monotonicity functional \\
$\mathcal{S}$ & Distributional scalar curvature measure \\
$\Energy_p(u)$ & $p$-energy: $\int |\nabla u|^p \dV$ \\
$\Cap_p(K)$ & $p$-capacity of compact set $K$ \\
$\lambda_1(L_\Sigma)$ & Principal eigenvalue of stability operator \\
\hline
\end{tabular}
\end{table}

\begin{table}[ht]
\centering
\caption{Weight and decay parameters.}
\begin{tabular}{l l l}
\hline
\textbf{Symbol} & \textbf{Range} & \textbf{Role} \\
\hline
$\tau$ & $> 1/2$ & Asymptotic flatness decay rate \\
$\delta$ & $(0, \tau - 1/2)$ & Weight for AF end (order $r^{-\delta}$) \\
$\beta$ & $(-1, 0)$ & Weight for cylindrical ends (order $e^{\beta t}$) \\
$\epsilon$ & $(0, \epsilon_0)$ & Smoothing parameter \\
$\kappa$ & $> 0$ & Surface gravity (blow-up rate: $f \sim \kappa^{-1} \ln s$) \\
\hline
\end{tabular}
\end{table}

\begin{table}[ht]
\centering
\caption{Jump and Interface Quantities.}
\begin{tabular}{l l l}
\hline
\textbf{Symbol} & \textbf{Definition} & \textbf{Context} \\
\hline
$[H]_{\bar{g}}$ & $H_{\text{outer}}^{\bar{g}} - H_{\text{inner}}^{\bar{g}}$ & Mean curvature jump in Jang metric $\bar{g}$ \\
$[H]_{\tilde{g}}$ & Jump in conformal metric $\tilde{g} = \phi^4 \bar{g}$ & Used in Miao smoothing \\
$\tr_\Sigma k$ & Trace of extrinsic curvature & Initial data quantity \\
$\theta^\pm$ & $H_\Sigma \pm \tr_\Sigma k$ & Null expansions \\
\hline
\end{tabular}
\end{table}

\subsection{Detailed Notation Dictionary for Metrics and Jumps}

To avoid ambiguity, we explicitly define the various metrics and jump quantities used in the interface analysis.

\begin{description}
    \item[Initial Data Metric ($g$):] The Riemannian metric on the initial slice $M$.
    \item[Jang Metric ($\bar{g}$):] The metric on the graph $\Sigma \subset M \times \mathbb{R}$, given by $\bar{g} = g + df \otimes df$. It is Lipschitz across the interface $\Sigma$.
    \item[Conformal Jang Metric ($\tilde{g}$):] The metric $\tilde{g} = \phi^4 \bar{g}$ used in the AMO flow. The conformal factor $\phi$ solves the Lichnerowicz equation.
    \item[Smoothed Metric ($g_\epsilon$):] The smooth family of metrics approximating $\tilde{g}$ (or $\bar{g}$) via the Miao--Piubello smoothing.
\end{description}

\textbf{Jump Conventions:}
For a quantity $Q$ discontinuous across a surface $\Sigma$, we define the jump $[Q] = Q^+ - Q^-$, where:
\begin{itemize}
    \item $Q^+$ is the limit from the \textbf{exterior} (asymptotically flat side).
    \item $Q^-$ is the limit from the \textbf{interior} (cylindrical/black hole side).
\end{itemize}
The unit normal $\nu$ points from interior to exterior.

\textbf{Key Identities:}
\begin{itemize}
    \item \textbf{Jang Jump:} $[H]_{\bar{g}} = \tr_\Sigma k$ (under the favorable jump hypothesis).
    \item \textbf{Conformal Jump:} $[H]_{\tilde{g}} = \phi^{-2} [H]_{\bar{g}} + 4\phi^{-3}[\partial_\nu \phi]$.
    \item \textbf{Transmission Condition:} If $[\partial_\nu \phi] = 0$, then $[H]_{\tilde{g}} = \phi^{-2} [H]_{\bar{g}}$.
\end{itemize}

\begin{remark}[Jump Relations]
The mean curvature jumps are related by the conformal factor $\phi$. If the transmission condition $[\partial_\nu \phi] = 0$ holds, then:
\[
[H]_{\tilde{g}} = \phi^{-2} [H]_{\bar{g}}.
\]
The fundamental identity relating the Jang metric jump to initial data is $[H]_{\bar{g}} = \tr_\Sigma k$ (Lemma~\ref{lem:TrappedMeanCurvatureJump}).
\end{remark}

\begin{table}[ht]
\centering
\caption{Key theorems and their roles.}
\small
\begin{tabular}{@{}lll@{}}
\hline
\textbf{Theorem} & \textbf{Content} & \textbf{Section} \\
\hline
\ref{thm:penroseinitial} & Primary result (outermost MOTS) & \S\ref{sec:intro} \\
\ref{thm:MainTheorem} & Conditional Penrose inequality & \S\ref{sec:penrose_conjecture} \\
\ref{thm:CompleteProof} & Consolidated proof & \S\ref{sec:Consolidated} \\
\ref{thm:MaxAreaTrapped} & Maximum area trapped surface & \S\ref{sec:Overview} \\
\ref{thm:AMOMonotonicity} & AMO level set monotonicity & \S\ref{sec:AMO} \\
\ref{thm:CompleteMeanCurvatureJump} & Mean curvature jump & \S\ref{sec:Jang} \\
\hline
\end{tabular}
\end{table}

\begin{table}[ht]
\centering
\caption{Sign conventions summary.}
\begin{tabular}{l l}
\hline
\textbf{Quantity} & \textbf{Convention} \\
\hline
Mean curvature $H$ & $H > 0$ for outward-bending surfaces \\
Null expansion $\theta^\pm$ & $\theta^\pm = H \pm \tr_\Sigma k$ \\
Trapped surface & $\theta^+ \le 0$ (outer trapped) \\
MOTS & $\theta^+ = 0$ (marginally outer trapped) \\
Stability operator $L_\Sigma$ & $\lambda_1 \ge 0$ for stable MOTS \\
Mean curvature jump $[H]$ & $[H] = H^+ - H^-$ (exterior minus interior) \\
Scalar curvature $R$ & Round sphere has $R > 0$ \\
\hline
\end{tabular}
\end{table}

\begin{table}[ht]
\centering
\caption{Geometric Quantities and Jumps.}
\begin{tabular}{l l l}
\hline
\textbf{Symbol} & \textbf{Description} & \textbf{Sign Convention} \\
\hline
$\theta^+$ & Outer null expansion & $\le 0$ for trapped \\
$\theta^-$ & Inner null expansion & $< 0$ for trapped \\
$\tr_\Sigma k$ & Trace of extrinsic curvature on $\Sigma$ & Favorable if $\ge 0$ \\
$[H]_{\bar{g}}$ & Mean curvature jump in Jang metric & $[H]_{\bar{g}} = \tr_\Sigma k$ \\
$[H]_{\tilde{g}}$ & Jump in conformal metric & $[H]_{\tilde{g}} = \phi^{-2}[H]_{\bar{g}}$ \\
$\lambda_1$ & Principal eigenvalue of stability op. & Stable if $\ge 0$ \\
\hline
\end{tabular}
\end{table}


\input{sec_32_worked_example_schwarzschild_initial_data}

\bibliographystyle{amsplain}
\bibliography{references}

\end{document}
